\scanpage{49}
% archive sheet number: 48

(\uline{20 suite})

$Z'\cap Z''$ est un ouvert \deleted{\ensuremath{\unsure{de|Z'}}} de $V\cap W=U$, d'ailleurs $Z$ est obtenu en recollant $Z'$ et $Z''$ suivant $I$ (ouvert commun). De \deleted{\unsure{ce|m|même}} $Z''\cap V$ est un ouvert de $V$, (mais peut-être pas de $Z''$) qui, avec son ouvert $Z'$, recouvre $V$, et $=$ en sens inverse.

% slanted
\margin{Trop compliqué, \unsure{d'abord|donc} établir le calcul de la structure de $T'$, la \unsure{compatib$^{\mathrm{té}}$|complétude} de la fonct\textsuperscript{n} de $e',e''$ aux \unsure{f$^*$|f}, de $f$ à \unsure{base|bas}, \unsure{sans|pour} un raisonn\textsuperscript{t}}

\uline{En résumé} : Pour que le Pb ait une solution, il f et s. \deleted{qu'il existe un ouvert $Z'$ de $V$} que $\exists Z'$ \deleted{est le plus gd} ouvert de $V$ qui induit \deleted{un} ouvert \deleted{sur} $Z'\cap W$ sur $W$, et $Z''$ \deleted{le $2$ plus gd} ouvert de $W$ qui induit un ouvert $Z''\cap V$ sur $V$, alors $V=\mathrm{Sup}(Z',Z''\cap V)$, $W=\mathrm{Sup}(Z'',Z'\cap W)$.

On a vu que c'est nécessaire. Pour voir que c'est suffisant, considérons le topos $\unsure{Z^s|Z|Z^*}$ qui est obtenu en recollant $Z'$ et $Z''$, suivant leur intersection $\unsure{I|Z|Z_1}$, d'où un morphisme $Z\to B$ qui est un plgmt (\uline{à vérifier}) et des morphismes $V\to\unsure{Z^s|Z|Z^*}$ et $W\to Z$ qui sont définis via $Z'\xrightarrow{\mathrm{Id}}Z'$ et $V\cap Z''\longrightarrow Z''$ \deleted{\ensuremath{\unsure{\times_W Z''|\times Z''}}} d'une part ($V=\mathrm{Sup}(Z',V\cap Z'')$) et via \ldots\ d'autre part, qui sont des plgmts (car composés avec $Z\to B$ on obtient les plgmts donnés), donc $V,W\subset Z\subset B$. On voit alors aussitôt que $Z$ satisfait à la condition

% slanted
\margin{construction à préciser}

\uline{NB} Il suffit en particulier que $Z'\cap Z''$ soit ouvert dans $Z'$ et dans $Z''$\ldots.

Dans le cas envisagé $B=\widehat R=X_{\mathrm{ann}}$, $V=X_{\mathrm{car}>0}$, $W=X_{\mathrm{car}\,0}$, on sait que $V$ est déjà un ouvert de $B=\widehat R$ (correspond au crible des objets $X\in\mathrm{Ob}\,R$ qui sont une car positive \deleted{\unsure{fixe|donnée}} \inserted{connue}, \unsure{s.les|sls|savoir}, un premier). NB ce crible \deleted{n'a pas des \unsure{limites produits|produits fibrés}, et} \inserted{a des produits fibrés}, \deleted{\unsure{voyons|mais|noyaux}} mais a des doubles flèches, (mais pas \unsure{des objet|de objet|d'objet} final) et $V\cap W$ est l'ouvert de $X_{\mathrm{car}\,0}\simeq\widehat{R_{\mathbb Q}}$ correspondant à \deleted{l'anneau nul} l'objet nul \inserted{(initial ?)} de $R_{\mathbb Q}$.

Donc \uuline{OK}

\margin{\uline{NB}\quad $X_{\mathrm{car}>0}\simeq\widehat{R_{\mathrm{car}>0}}$, $R_{\mathrm{car}>0}$ sous-cat \unsure{apl.|pl.} de $R$ des schémas ayant une car, $=$ somme amalgamée des $R_{\mathbb Z/p\mathbb Z}$ (\unsure{pr finie|p fini}) sur $\varnothing_R\longrightarrow R_{\mathbb Z/p\mathbb Z}$.}
